\documentclass[letterpaper, 12pt]{article}

% ============================================================
% PAQUETES
% ============================================================
\usepackage[utf8]{inputenc}
\usepackage[spanish]{babel}
\usepackage{amsmath, amssymb, amsfonts}
\usepackage{geometry}
\usepackage{fancyhdr}
\usepackage{enumitem}
\usepackage{graphicx}
\usepackage{hyperref}
\usepackage{multicol}
\usepackage{parskip}
\usepackage{xcolor}
\usepackage{tcolorbox}
\usepackage{eso-pic}
\usepackage{transparent}
\usepackage{cancel}
\tcbuselibrary{skins}

% ============================================================
% GEOMETRÍA
% ============================================================
\geometry{letterpaper, margin=1.5cm, top=3cm, bottom=2.5cm, headheight=2cm}

% ============================================================
% NOTACIÓN ESPAÑOLA
% ============================================================
\let\sin\relax
\DeclareMathOperator{\sin}{sen}

% ============================================================
% COLORES INSTITUCIONALES
% ============================================================
\definecolor{celesteSuave}{HTML}{F0F7FF}
\definecolor{azulFuerte}{HTML}{1A5276}
\definecolor{enlace}{HTML}{1A5276}
\definecolor{verdeOk}{HTML}{27AE60}
\definecolor{rojoAlerta}{HTML}{C0392B}

% ============================================================
% INFORMACIÓN DE LA CLASE
% ============================================================
\newcommand{\logoup}{./images/logo_up.png}
\newcommand{\logousach}{./images/logo_usach.png}
\newcommand{\logofondo}{./images/logo_up.png}

\newcommand{\institucion}{Usach Premium}
\newcommand{\curso}{Cálculo I}
\newcommand{\guiasnumero}{Guía 2 — Clase en Vivo (Uso Docente)}
\newcommand{\semestre}{1er. Semestre de 2026}
\newcommand{\fraseportada}{"Por y para estudiantes"}

% ============================================================
% HYPERREF
% ============================================================
\hypersetup{
    colorlinks=true,
    linkcolor=enlace,
    urlcolor=enlace,
    citecolor=enlace,
    pdfborder={0 0 0},
}

% ============================================================
% ENCABEZADO Y PIE DE PÁGINA
% ============================================================
\pagestyle{fancy}
\fancyhf{}
\fancyhead[L]{%
    \includegraphics[height=1.2cm]{\logoup}%
}
\fancyhead[R]{%
    \begin{minipage}[b]{0.50\textwidth}
        \raggedleft
        \fontsize{9pt}{11pt}\selectfont
        \textbf{\color{azulFuerte}\institucion} \\
        \curso\ — Clase en Vivo \\
        \textit{\color{gray}Documento de uso docente}
    \end{minipage}\hspace{0.1cm}%
    \includegraphics[height=1.2cm]{\logousach}%
}
\fancyfoot[L]{\fontsize{9pt}{11pt}\selectfont \textit{\fraseportada}}
\fancyfoot[R]{\fontsize{9pt}{11pt}\selectfont Página \thepage}
\renewcommand{\headrulewidth}{0.4pt}
\renewcommand{\footrulewidth}{0.4pt}

\fancypagestyle{plain}{
    \fancyhf{}
    \renewcommand{\headrulewidth}{0pt}
    \renewcommand{\footrulewidth}{0pt}
}

% ============================================================
% MARCA DE AGUA
% ============================================================
\newcommand{\MarcaAgua}{
    \AddToShipoutPicture{
        \AtPageCenter{
            \makebox(0,0){%
                \transparent{0.05}%
                \includegraphics[width=0.7\textwidth]{\logofondo}%
            }
        }
    }
}

% ============================================================
% CAJAS TCOLORBOX
% ============================================================
\newtcolorbox{materiabox}[1]{
    colback=celesteSuave,
    colframe=azulFuerte,
    fonttitle=\bfseries,
    coltitle=white,
    title=#1,
    arc=8pt,
    boxrule=1.2pt,
    enhanced,
    drop shadow={black!5}
}

\newtcolorbox{solucionbox}[1]{
    colback=white,
    colframe=verdeOk,
    fonttitle=\bfseries,
    coltitle=white,
    title=#1,
    arc=6pt,
    boxrule=1pt,
    enhanced,
    drop shadow={black!5}
}

\newtcolorbox{desafiobox}[1]{
    colback=white,
    colframe=azulFuerte,
    fonttitle=\bfseries,
    coltitle=white,
    title=#1,
    arc=8pt,
    boxrule=1.2pt,
    enhanced,
    drop shadow={black!5}
}

% ============================================================
\begin{document}
\thispagestyle{plain}
\MarcaAgua

% ============================================================
% PORTADA
% ============================================================
\AddToShipoutPicture*{%
    \put(54, 700){\includegraphics[width=2cm]{\logoup}}
    \put(420, 706){\includegraphics[width=5cm]{\logousach}}
}

\vspace*{0.5cm}

\begin{center}
    \color{azulFuerte}
    {\huge \textbf{CLASE EN VIVO — EJERCICIOS}} \\[0.5cm]
    {\Huge \textbf{\curso}} \\[0.3cm]
    {\Large Guía 2 — Nivelación} \\[0.3cm]
    {\large\color{rojoAlerta}\textbf{DOCUMENTO DE USO DOCENTE — NO ENTREGAR}}\\[1.2cm]

    \begin{tcolorbox}[colback=celesteSuave, colframe=azulFuerte, arc=10pt,
        width=0.85\textwidth, center, boxrule=1.5pt]
        \centering
        \vspace{0.2cm}
        {\Large \textbf{Contenidos de la Clase}}
        \vspace{0.3cm}
        \color{black}
        \begin{itemize}[leftmargin=3cm]
            \item[\color{azulFuerte}$\Rightarrow$] Materia: Ecuaciones Lineales y Cuadráticas
            \item[\color{azulFuerte}$\Rightarrow$] Materia: Razones Trigonométricas
            \item[\color{azulFuerte}$\Rightarrow$] Ejercicios con Desarrollo Completo
            \item[\color{azulFuerte}$\Rightarrow$] Desafíos N1 y N2 Resueltos
        \end{itemize}
        \vspace{0.2cm}
    \end{tcolorbox}
\end{center}

\vspace{1.5cm}

\begin{center}
    {\Large \textbf{\semestre}} \\[0.5cm]
    {\large \textbf{Byron Caices y Luis Vásquez}}
\end{center}

\pagenumbering{arabic}
\newpage

% ============================================================
% MATERIA: ECUACIONES
% ============================================================
\section*{Resumen de Materia}

\begin{materiabox}{ECUACIONES LINEALES}
Una \textbf{ecuación lineal} es de la forma $ax + b = 0$, con $a \neq 0$.

\textbf{Estrategia de resolución:}
\begin{enumerate}[leftmargin=*, itemsep=0.15cm]
    \item Eliminar paréntesis (propiedad distributiva).
    \item Eliminar denominadores multiplicando por el MCM.
    \item Agrupar términos con $x$ a un lado y constantes al otro.
    \item Despejar $x$ dividiendo por el coeficiente.
\end{enumerate}
\end{materiabox}

\vspace{0.4cm}

\begin{materiabox}{ECUACIONES DE SEGUNDO GRADO}
Forma general: $ax^2 + bx + c = 0$, con $a \neq 0$.

\textbf{Métodos de resolución:}
\begin{itemize}[leftmargin=*, itemsep=0.15cm]
    \item \textbf{Factorización:} Buscar dos números que multiplicados den $c$ y sumados den $b$.
    \item \textbf{Fórmula general:}
    \[
        x = \frac{-b \pm \sqrt{b^2 - 4ac}}{2a}
    \]
    \item \textbf{Discriminante} $\Delta = b^2 - 4ac$:
    \begin{itemize}
        \item $\Delta > 0$: dos soluciones reales distintas.
        \item $\Delta = 0$: una solución real doble.
        \item $\Delta < 0$: no hay soluciones reales.
    \end{itemize}
\end{itemize}
\end{materiabox}

\vspace{0.4cm}

\begin{materiabox}{RAZONES TRIGONOMÉTRICAS}
En un triángulo rectángulo con ángulo $\alpha$:
\begin{multicols}{2}
\begin{itemize}[leftmargin=*, itemsep=0.15cm]
    \item $\sin \alpha = \dfrac{\text{Cateto Opuesto}}{\text{Hipotenusa}}$
    \item $\cos \alpha = \dfrac{\text{Cateto Adyacente}}{\text{Hipotenusa}}$
    \item $\tan \alpha = \dfrac{\text{Cateto Opuesto}}{\text{Cateto Adyacente}}$
    \item $\csc \alpha = \dfrac{1}{\sin \alpha}$
    \item $\sec \alpha = \dfrac{1}{\cos \alpha}$
    \item $\cot \alpha = \dfrac{1}{\tan \alpha}$
\end{itemize}
\end{multicols}

\textbf{Identidades Fundamentales:}
\begin{itemize}[leftmargin=*, itemsep=0.15cm]
    \item \textbf{Pitagórica:} $\sin^2\alpha + \cos^2\alpha = 1$
    \item \textbf{Derivada 1:} $1 + \tan^2\alpha = \sec^2\alpha$
    \item \textbf{Derivada 2:} $1 + \cot^2\alpha = \csc^2\alpha$
\end{itemize}

\vspace{0.2cm}
\textbf{Valores Notables:}
\begin{center}
\renewcommand{\arraystretch}{1.4}
\begin{tabular}{|c|c|c|c|c|}
\hline
 & $\mathbf{30^\circ}$ & $\mathbf{45^\circ}$ & $\mathbf{60^\circ}$ & $\mathbf{90^\circ}$ \\ \hline
$\sin$ & $1/2$ & $\sqrt{2}/2$ & $\sqrt{3}/2$ & $1$ \\ \hline
$\cos$ & $\sqrt{3}/2$ & $\sqrt{2}/2$ & $1/2$ & $0$ \\ \hline
$\tan$ & $\sqrt{3}/3$ & $1$ & $\sqrt{3}$ & $\infty$ \\ \hline
\end{tabular}
\end{center}
\end{materiabox}

\newpage

% ============================================================
% SECCIÓN 1: ECUACIONES LINEALES (5 ejercicios con desarrollo)
% ============================================================
\section*{Sección 1: Ecuaciones Lineales — Ejercicios con Desarrollo}

% --- Ejercicio 1 ---
\textbf{Ejercicio 1.} Resuelve: $3(x - 2) = 12$

\begin{solucionbox}{Desarrollo}
\begin{align*}
    3(x - 2) &= 12 \\
    3x - 6 &= 12 && \text{(Distributiva)} \\
    3x &= 12 + 6 \\
    3x &= 18 \\
    x &= \frac{18}{3} \\[4pt]
    &\boxed{x = 6}
\end{align*}
\end{solucionbox}

\vspace{0.4cm}

% --- Ejercicio 2 ---
\textbf{Ejercicio 2.} Resuelve: $4(2x - 1) - 3(x + 2) = 10$

\begin{solucionbox}{Desarrollo}
\begin{align*}
    4(2x - 1) - 3(x + 2) &= 10 \\
    8x - 4 - 3x - 6 &= 10 && \text{(Distributiva)} \\
    5x - 10 &= 10 && \text{(Reducir términos semejantes)} \\
    5x &= 20 \\
    x &= \frac{20}{5} \\[4pt]
    &\boxed{x = 4}
\end{align*}
\end{solucionbox}

\vspace{0.4cm}

% --- Ejercicio 3 ---
\textbf{Ejercicio 3.} Resuelve: $\displaystyle\frac{x}{2} + \frac{x}{3} = 5$

\begin{solucionbox}{Desarrollo}
El MCM de $2$ y $3$ es $6$. Multiplicamos toda la ecuación por $6$:
\begin{align*}
    \frac{x}{2} + \frac{x}{3} &= 5 \\[6pt]
    6 \cdot \frac{x}{2} + 6 \cdot \frac{x}{3} &= 6 \cdot 5 \\[6pt]
    3x + 2x &= 30 \\
    5x &= 30 \\
    x &= \frac{30}{5} \\[4pt]
    &\boxed{x = 6}
\end{align*}
\end{solucionbox}

\vspace{0.4cm}

% --- Ejercicio 4 ---
\textbf{Ejercicio 4.} Resuelve: $\displaystyle\frac{2x - 1}{3} = \frac{x + 4}{2}$

\begin{solucionbox}{Desarrollo}
Multiplicamos cruzado (o por el MCM $= 6$):
\begin{align*}
    \frac{2x - 1}{3} &= \frac{x + 4}{2} \\[6pt]
    2(2x - 1) &= 3(x + 4) && \text{(Multiplicación cruzada)} \\
    4x - 2 &= 3x + 12 \\
    4x - 3x &= 12 + 2 \\
    x &= 14 \\[4pt]
    &\boxed{x = 14}
\end{align*}
\end{solucionbox}

\vspace{0.4cm}

% --- Ejercicio 5 ---
\textbf{Ejercicio 5.} Resuelve: $\displaystyle\frac{x - 2}{4} - \frac{x + 1}{3} = -1$

\begin{solucionbox}{Desarrollo}
El MCM de $4$ y $3$ es $12$. Multiplicamos toda la ecuación por $12$:
\begin{align*}
    \frac{x - 2}{4} - \frac{x + 1}{3} &= -1 \\[6pt]
    12 \cdot \frac{x-2}{4} - 12 \cdot \frac{x+1}{3} &= 12 \cdot (-1) \\[6pt]
    3(x - 2) - 4(x + 1) &= -12 \\
    3x - 6 - 4x - 4 &= -12 \\
    -x - 10 &= -12 \\
    -x &= -12 + 10 \\
    -x &= -2 \\
    x &= 2 \\[4pt]
    &\boxed{x = 2}
\end{align*}
\end{solucionbox}

\newpage

% ============================================================
% SECCIÓN 2: ECUACIONES DE SEGUNDO GRADO (5 ejercicios con desarrollo)
% ============================================================
\section*{Sección 2: Ecuaciones de Segundo Grado — Ejercicios con Desarrollo}

% --- Ejercicio 6 ---
\textbf{Ejercicio 6.} Resuelve: $x^2 - 9 = 0$

\begin{solucionbox}{Desarrollo — Diferencia de Cuadrados}
\begin{align*}
    x^2 - 9 &= 0 \\
    x^2 - 3^2 &= 0 && \text{(Diferencia de cuadrados: } a^2 - b^2 = (a-b)(a+b)\text{)} \\
    (x - 3)(x + 3) &= 0
\end{align*}
Igualando cada factor a cero:
\[
    x - 3 = 0 \implies x = 3 \qquad \text{o} \qquad x + 3 = 0 \implies x = -3
\]
\[\boxed{x = 3 \quad \text{o} \quad x = -3}\]
\end{solucionbox}

\vspace{0.4cm}

% --- Ejercicio 7 ---
\textbf{Ejercicio 7.} Resuelve: $x^2 - 5x + 6 = 0$

\begin{solucionbox}{Desarrollo — Factorización}
Buscamos dos números que multiplicados den $+6$ y sumados den $-5$: esos son $-2$ y $-3$.
\begin{align*}
    x^2 - 5x + 6 &= 0 \\
    (x - 2)(x - 3) &= 0
\end{align*}
Igualando cada factor a cero:
\[
    x - 2 = 0 \implies x = 2 \qquad \text{o} \qquad x - 3 = 0 \implies x = 3
\]
\[\boxed{x = 2 \quad \text{o} \quad x = 3}\]
\end{solucionbox}

\vspace{0.4cm}

% --- Ejercicio 8 ---
\textbf{Ejercicio 8.} Resuelve: $x^2 - 4x - 5 = 0$

\begin{solucionbox}{Desarrollo — Factorización}
Buscamos dos números que multiplicados den $-5$ y sumados den $-4$: esos son $-5$ y $+1$.
\begin{align*}
    x^2 - 4x - 5 &= 0 \\
    (x - 5)(x + 1) &= 0
\end{align*}
Igualando cada factor a cero:
\[
    x - 5 = 0 \implies x = 5 \qquad \text{o} \qquad x + 1 = 0 \implies x = -1
\]
\[\boxed{x = 5 \quad \text{o} \quad x = -1}\]
\end{solucionbox}

\vspace{0.4cm}

% --- Ejercicio 9 ---
\textbf{Ejercicio 9.} Resuelve: $3x^2 - 5x + 2 = 0$

\begin{solucionbox}{Desarrollo — Fórmula General}
Con $a = 3$, $b = -5$, $c = 2$:
\begin{align*}
    \Delta &= b^2 - 4ac = (-5)^2 - 4(3)(2) = 25 - 24 = 1 \\[6pt]
    x &= \frac{-b \pm \sqrt{\Delta}}{2a} = \frac{5 \pm \sqrt{1}}{6} = \frac{5 \pm 1}{6}
\end{align*}
\[
    x_1 = \frac{5 + 1}{6} = \frac{6}{6} = 1 \qquad \text{y} \qquad x_2 = \frac{5 - 1}{6} = \frac{4}{6} = \frac{2}{3}
\]
\[\boxed{x = 1 \quad \text{o} \quad x = \frac{2}{3}}\]

\textit{Verificación por factorización:} $3x^2 - 5x + 2 = (3x - 2)(x - 1) = 0$. \checkmark
\end{solucionbox}

\vspace{0.4cm}

% --- Ejercicio 10 ---
\textbf{Ejercicio 10.} Resuelve: $2x^2 + 7x - 4 = 0$

\begin{solucionbox}{Desarrollo — Fórmula General}
Con $a = 2$, $b = 7$, $c = -4$:
\begin{align*}
    \Delta &= b^2 - 4ac = 7^2 - 4(2)(-4) = 49 + 32 = 81 \\[6pt]
    x &= \frac{-b \pm \sqrt{\Delta}}{2a} = \frac{-7 \pm \sqrt{81}}{4} = \frac{-7 \pm 9}{4}
\end{align*}
\[
    x_1 = \frac{-7 + 9}{4} = \frac{2}{4} = \frac{1}{2} \qquad \text{y} \qquad x_2 = \frac{-7 - 9}{4} = \frac{-16}{4} = -4
\]
\[\boxed{x = \frac{1}{2} \quad \text{o} \quad x = -4}\]
\end{solucionbox}

\newpage

% ============================================================
% SECCIÓN 3: RAZONES E IDENTIDADES TRIGONOMÉTRICAS (5 ejercicios)
% ============================================================
\section*{Sección 3: Razones e Identidades Trigonométricas — Ejercicios con Desarrollo}

\textit{Nota: en esta clase trabajamos solo con razones trigonométricas en el triángulo rectángulo, no con funciones trigonométricas.}

\vspace{0.4cm}

% --- Ejercicio 11 ---
\textbf{Ejercicio 11.} En un triángulo rectángulo, si $\sin(\alpha) = \dfrac{3}{5}$, calcula $\cos(\alpha)$ y $\tan(\alpha)$.

\begin{solucionbox}{Desarrollo}
Sabemos que $\sin(\alpha) = \dfrac{\text{opuesto}}{\text{hipotenusa}} = \dfrac{3}{5}$, por lo que:
\begin{itemize}[leftmargin=*]
    \item Cateto opuesto $= 3$, Hipotenusa $= 5$.
\end{itemize}
Usamos el Teorema de Pitágoras para el cateto adyacente:
\begin{align*}
    b^2 + 3^2 &= 5^2 \\
    b^2 &= 25 - 9 = 16 \\
    b &= 4
\end{align*}
Entonces:
\[
    \cos(\alpha) = \frac{\text{adyacente}}{\text{hipotenusa}} = \frac{4}{5} \qquad \text{y} \qquad \tan(\alpha) = \frac{\text{opuesto}}{\text{adyacente}} = \frac{3}{4}
\]
\[\boxed{\cos(\alpha) = \frac{4}{5}, \qquad \tan(\alpha) = \frac{3}{4}}\]
\end{solucionbox}

\vspace{0.4cm}

% --- Ejercicio 12 ---
\textbf{Ejercicio 12.} Simplifica: $\sin^2(x) + \cos^2(x) + \tan^2(x)$

\begin{solucionbox}{Desarrollo}
Por la identidad pitagórica fundamental: $\sin^2(x) + \cos^2(x) = 1$.

Sustituyendo:
\begin{align*}
    \sin^2(x) + \cos^2(x) + \tan^2(x) &= 1 + \tan^2(x)
\end{align*}
Aplicando la identidad derivada $1 + \tan^2(x) = \sec^2(x)$:
\[\boxed{\sin^2(x) + \cos^2(x) + \tan^2(x) = \sec^2(x)}\]
\end{solucionbox}

\vspace{0.4cm}

% --- Ejercicio 13 ---
\textbf{Ejercicio 13.} Simplifica: $(1 - \sin^2(x)) \cdot \sec^2(x)$

\begin{solucionbox}{Desarrollo}
Paso 1: Aplicamos la identidad pitagórica $\sin^2(x) + \cos^2(x) = 1$, de donde $1 - \sin^2(x) = \cos^2(x)$.

Paso 2: Recordamos que $\sec(x) = \dfrac{1}{\cos(x)}$, por lo tanto $\sec^2(x) = \dfrac{1}{\cos^2(x)}$.

\begin{align*}
    (1 - \sin^2(x)) \cdot \sec^2(x) &= \cos^2(x) \cdot \frac{1}{\cos^2(x)} = 1
\end{align*}
\[\boxed{(1 - \sin^2(x)) \cdot \sec^2(x) = 1}\]
\end{solucionbox}

\vspace{0.4cm}

% --- Ejercicio 14 ---
\textbf{Ejercicio 14.} Calcula el valor numérico exacto de $\sin(30^\circ) + \cos(60^\circ)$.

\begin{solucionbox}{Desarrollo}
Usando la tabla de valores notables:
\[
    \sin(30^\circ) = \frac{1}{2} \qquad \text{y} \qquad \cos(60^\circ) = \frac{1}{2}
\]
Sumando:
\[
    \sin(30^\circ) + \cos(60^\circ) = \frac{1}{2} + \frac{1}{2} = 1
\]

\textit{Observación: $30^\circ$ y $60^\circ$ son complementarios, y $\sin(\alpha) = \cos(90^\circ - \alpha)$.}

\[\boxed{\sin(30^\circ) + \cos(60^\circ) = 1}\]
\end{solucionbox}

\vspace{0.4cm}

% --- Ejercicio 15 ---
\textbf{Ejercicio 15.} Simplifica: $\cos(x)(\sec(x) - \cos(x))$

\begin{solucionbox}{Desarrollo}
Distribuimos $\cos(x)$:
\begin{align*}
    \cos(x)\bigl(\sec(x) - \cos(x)\bigr) &= \cos(x) \cdot \sec(x) - \cos^2(x)
\end{align*}
Sabemos que $\cos(x) \cdot \sec(x) = \cos(x) \cdot \dfrac{1}{\cos(x)} = 1$. Luego:
\begin{align*}
    &= 1 - \cos^2(x)
\end{align*}
Aplicando la identidad pitagórica $1 - \cos^2(x) = \sin^2(x)$:
\[\boxed{\cos(x)(\sec(x) - \cos(x)) = \sin^2(x)}\]
\end{solucionbox}

\newpage

% ============================================================
% SECCIÓN 4: DESAFÍOS N1 — EXPRESIONES ALGEBRAICAS (con desarrollo)
% ============================================================
\section*{Desafíos N1: Expresiones Algebraicas — Desarrollo Completo}

\begin{desafiobox}{DESAFÍO 41 — Simplificación Racional Compleja}
Simplifica la siguiente expresión a su forma más irreductible:
\[
    \frac{x^4 - 16}{x^3 + 2x^2 + 4x + 8} \div \frac{x^2 - 4x + 4}{x^2 - 4}
\]
\end{desafiobox}

\begin{solucionbox}{Desarrollo Completo}

\textbf{Paso 1: Factorizar cada polinomio.}

\begin{itemize}[leftmargin=*, itemsep=0.3cm]
    \item $x^4 - 16 = (x^2)^2 - 4^2 = (x^2 - 4)(x^2 + 4) = (x-2)(x+2)(x^2+4)$

    \item $x^3 + 2x^2 + 4x + 8$: agrupamos
    \[
        = x^2(x + 2) + 4(x + 2) = (x^2 + 4)(x + 2)
    \]

    \item $x^2 - 4x + 4 = (x - 2)^2$

    \item $x^2 - 4 = (x - 2)(x + 2)$
\end{itemize}

\textbf{Paso 2: Reescribir la división como multiplicación.}
\[
    \frac{x^4 - 16}{x^3 + 2x^2 + 4x + 8} \cdot \frac{x^2 - 4}{x^2 - 4x + 4}
\]

\textbf{Paso 3: Sustituir las factorizaciones.}
\[
    = \frac{(x-2)(x+2)(x^2+4)}{(x^2+4)(x+2)} \cdot \frac{(x-2)(x+2)}{(x-2)^2}
\]

\textbf{Paso 4: Cancelar factores comunes.}
\[
    = \frac{\cancel{(x-2)}\cancel{(x+2)}\cancel{(x^2+4)}}{\cancel{(x^2+4)}\cancel{(x+2)}} \cdot \frac{\cancel{(x-2)}(x+2)}{(x-2)^{\cancel{2}}}
\]
\[
    = \frac{x + 2}{1} = x + 2
\]

\[\boxed{\frac{x^4 - 16}{x^3 + 2x^2 + 4x + 8} \div \frac{x^2 - 4x + 4}{x^2 - 4} = x + 2}\]

\textit{Restricciones de dominio:} $x \neq 2$, $x \neq -2$.
\end{solucionbox}

\newpage

\begin{desafiobox}{DESAFÍO 42 — Factorización Avanzada}
Factoriza completamente:
\[
    x^6 - 64y^6
\]
\end{desafiobox}

\begin{solucionbox}{Desarrollo Completo}

\textbf{Paso 1: Reconocer como diferencia de cuadrados.}
\[
    x^6 - 64y^6 = (x^3)^2 - (8y^3)^2
\]

\textit{Nota:} $(2y)^6 = 64y^6$, y $(2y)^3 = 8y^3$.

Aplicando $a^2 - b^2 = (a - b)(a + b)$:
\[
    = (x^3 - 8y^3)(x^3 + 8y^3)
\]

\textbf{Paso 2: Factorizar cada factor como suma/diferencia de cubos.}

Recordamos:
\begin{itemize}
    \item $a^3 - b^3 = (a - b)(a^2 + ab + b^2)$
    \item $a^3 + b^3 = (a + b)(a^2 - ab + b^2)$
\end{itemize}

\textbf{Diferencia de cubos:} $x^3 - 8y^3 = x^3 - (2y)^3$
\[
    = (x - 2y)(x^2 + 2xy + 4y^2)
\]

\textbf{Suma de cubos:} $x^3 + 8y^3 = x^3 + (2y)^3$
\[
    = (x + 2y)(x^2 - 2xy + 4y^2)
\]

\textbf{Paso 3: Resultado final.}
\[\boxed{x^6 - 64y^6 = (x - 2y)(x^2 + 2xy + 4y^2)(x + 2y)(x^2 - 2xy + 4y^2)}\]
\end{solucionbox}

\vspace{0.6cm}

\begin{desafiobox}{DESAFÍO 43 — Fracciones Algebraicas}
Simplifica la siguiente expresión compuesta, indicando restricciones de dominio:
\[
    \frac{\dfrac{1}{x} - \dfrac{1}{x+h}}{h}
\]
\end{desafiobox}

\begin{solucionbox}{Desarrollo Completo}

\textbf{Paso 1: Operar la resta de fracciones del numerador.}
\[
    \frac{1}{x} - \frac{1}{x+h} = \frac{(x+h) - x}{x(x+h)} = \frac{h}{x(x+h)}
\]

\textbf{Paso 2: Dividir por $h$.}
\[
    \frac{\dfrac{h}{x(x+h)}}{h} = \frac{h}{x(x+h)} \cdot \frac{1}{h} = \frac{\cancel{h}}{x(x+h) \cdot \cancel{h}} = \frac{1}{x(x+h)}
\]

\[\boxed{\frac{\dfrac{1}{x} - \dfrac{1}{x+h}}{h} = \frac{1}{x(x+h)}}\]

\textit{Restricciones de dominio:} $x \neq 0$, \ $x \neq -h$, \ $h \neq 0$.

\vspace{0.2cm}
\textit{Observación:} Esta expresión es el \textbf{cuociente de diferencias} de $f(x) = \dfrac{1}{x}$, que es la base del concepto de derivada en Cálculo.
\end{solucionbox}

\newpage

% ============================================================
% SECCIÓN 5: DESAFÍOS N2 — ECUACIONES Y TRIGONOMETRÍA (con desarrollo)
% ============================================================
\section*{Desafíos N2: Ecuaciones y Trigonometría — Desarrollo Completo}

\begin{desafiobox}{DESAFÍO 44 — Ecuación Paramétrica}
Resuelve para $x$, con $a, b$ constantes reales positivas y $a \neq -b$:
\[
    a^2\,x - b^2 = a^2 - b^2\,x
\]
\end{desafiobox}

\begin{solucionbox}{Desarrollo Completo}

\textbf{Paso 1: Agrupar los términos con $x$ a un lado.}
\begin{align*}
    a^2\,x - b^2 &= a^2 - b^2\,x \\
    a^2\,x + b^2\,x &= a^2 + b^2 && \text{(Sumamos } b^2 x \text{ y } b^2 \text{ a ambos lados)}
\end{align*}

\textbf{Paso 2: Factorizar $x$.}
\begin{align*}
    x(a^2 + b^2) &= a^2 + b^2
\end{align*}

\textbf{Paso 3: Despejar $x$.}

Como $a$ y $b$ son reales positivos, $a^2 + b^2 > 0$, luego podemos dividir:
\[
    x = \frac{a^2 + b^2}{a^2 + b^2} = 1
\]

\[\boxed{x = 1}\]

\textit{Nota: la condición $a \neq -b$ garantiza que $a + b \neq 0$, pero la clave aquí es que $a^2 + b^2 > 0$ para todo $a, b$ positivos, lo que siempre permite dividir.}
\end{solucionbox}

\vspace{0.6cm}

\begin{desafiobox}{DESAFÍO 45 — Identidad Trigonométrica Compleja}
Demuestra:
\[
    \frac{\sin(x) + \sin(2x)}{1 + \cos(x) + \cos(2x)} = \tan(x)
\]
\end{desafiobox}

\begin{solucionbox}{Desarrollo Completo}

\textbf{Paso 1: Reescribir usando identidades de ángulo doble.}

Recordamos:
\begin{itemize}
    \item $\sin(2x) = 2\sin(x)\cos(x)$
    \item $\cos(2x) = 2\cos^2(x) - 1$
\end{itemize}

\textbf{Paso 2: Trabajar el numerador.}
\begin{align*}
    \sin(x) + \sin(2x) &= \sin(x) + 2\sin(x)\cos(x) \\
    &= \sin(x)\bigl(1 + 2\cos(x)\bigr)
\end{align*}

\textbf{Paso 3: Trabajar el denominador.}
\begin{align*}
    1 + \cos(x) + \cos(2x) &= 1 + \cos(x) + 2\cos^2(x) - 1 \\
    &= \cos(x) + 2\cos^2(x) \\
    &= \cos(x)\bigl(1 + 2\cos(x)\bigr)
\end{align*}

\textbf{Paso 4: Dividir.}
\[
    \frac{\sin(x)\cancel{\bigl(1 + 2\cos(x)\bigr)}}{\cos(x)\cancel{\bigl(1 + 2\cos(x)\bigr)}} = \frac{\sin(x)}{\cos(x)} = \tan(x)
\]

\[\boxed{\frac{\sin(x) + \sin(2x)}{1 + \cos(x) + \cos(2x)} = \tan(x) \qquad \blacksquare}\]
\end{solucionbox}

\newpage

\begin{desafiobox}{DESAFÍO 46 — Problema de Elevación de Múltiples Etapas}
Un topógrafo nota que el ángulo de elevación a la cima de una torre es de $30^\circ$. Al avanzar 20 metros hacia la torre, el ángulo de elevación cambia a $60^\circ$. ¿Cuál es la altura exacta de la torre?
\end{desafiobox}

\begin{solucionbox}{Desarrollo Completo}

\textbf{Paso 1: Definir variables.}

Sea $h$ la altura de la torre y $d$ la distancia horizontal desde el \textbf{segundo punto} de observación hasta la base de la torre.

\textbf{Paso 2: Plantear ecuaciones con las razones trigonométricas.}

\textit{Desde el segundo punto} (más cercano), el ángulo es $60^\circ$:
\[
    \tan(60^\circ) = \frac{h}{d} \implies h = d\sqrt{3} \qquad \cdots\ (1)
\]

\textit{Desde el primer punto} (20 m más lejos), el ángulo es $30^\circ$:
\[
    \tan(30^\circ) = \frac{h}{d + 20} \implies h = \frac{d + 20}{\sqrt{3}} \qquad \cdots\ (2)
\]

\textbf{Paso 3: Igualar las expresiones de $h$.}
\begin{align*}
    d\sqrt{3} &= \frac{d + 20}{\sqrt{3}}
\end{align*}
Multiplicamos ambos lados por $\sqrt{3}$:
\begin{align*}
    3d &= d + 20 \\
    2d &= 20 \\
    d &= 10 \text{ m}
\end{align*}

\textbf{Paso 4: Obtener la altura.}
\[
    h = d\sqrt{3} = 10\sqrt{3} \text{ m}
\]

\[\boxed{h = 10\sqrt{3} \approx 17{,}32 \text{ metros}}\]
\end{solucionbox}

\vspace{1.5cm}

\begin{center}
    {\color{azulFuerte}\hrule height 0.5pt}
    \vspace{0.5cm}
    \small\textbf{Usach Premium} — Documento de uso docente para clase en vivo. \\[0.2cm]
    \textit{\fraseportada}
\end{center}

\end{document}
